\section{Other modes, and how far the results reach}
\label{sec:modes}

\subsection{Direct succession}

Live overlap is one way for an episode to pass from source to target. Another
is immediate. The last event of the source directly generates the first event
of the target, with no stage at which both are live. We keep this mode
separate.

\begin{definition}[Direct succession]
% claim: TTD-01
\label{def:direct}
A target event is a \emph{direct successor} of a source event when the pair
passes the computed direct-succession check (Appendix~\ref{app:index}). A
direct successor is generated by that single, necessary source event and
carries the source's whole eight-field frontier, advanced by one strict step.
In addition, the source event is used once only, nothing is stored or replayed
in between, there is no alternative parent, and route and ownership are
transferred correctly.
\end{definition}

\begin{proposition}[Validity of direct succession]
% claim: TTD-01
\label{prop:direct}
Every model in the computed direct mode yields a valid direct-successor
certificate.
\end{proposition}

The conditions listed in Definition~\ref{def:direct} are necessary for passing
the check. We do not prove that they are sufficient, so they summarize the
check and do not characterize it.

\begin{proposition}[What direct succession excludes]
% claim: TTD-03
\label{prop:direct-exclusions}
The direct criterion excludes four mechanisms: detached copies, replay from an
archive, reconstruction through storage, and gaps. Reachability by data alone
does not make a direct successor. Conversely, a separate data path can coexist
with genuine direct succession.
\end{proposition}

So the literal claim that direct succession excludes every data path is false,
and we do not make it.

\begin{proposition}[Two different modes]
% claim: TTD-02
\label{prop:modes}
Strict overlap continuation and direct succession are incomparable. Each has
instances that are not instances of the other. In fact the two modes are
disjoint on every model, since overlap needs the source and target events in
one stage and direct succession needs them in adjacent stages. Both embed in a
broader, separately named relation of continuity paths, in which compatible
mixed paths compose and closed paths are excluded.
\end{proposition}

The broader relation never replaces strict continuation. When we say
continuation without qualification we mean the strict, live-overlap kind.

\subsection{Unbounded maintenance}

Theorem~\ref{thm:autonomy} covers every finite horizon. To go further we need a
policy that never runs out.

\begin{theorem}[Unbounded maintenance on a finite factor]
% claim: TTS-04
\label{thm:unbounded}
On the chain model, take the target's computed policy on a finite factor of its
state, with unbounded event registers. Assume that the policy's inputs are
ready, that the initial factor lies in the policy's computed greatest viable
set, and that the initial balance covers the computed reserve. The policy's
responses are owned by the target and normalized, and its positive outcomes
advance and credit at least what they debit. Then for every
admissible stream of disturbances and every stream of outcomes of positive
probability, the resulting infinite run has the following properties at every
step:
\begin{itemize}
\item it stays in the policy's greatest viable set and realizes the
conditions, with (4) to (7) in simplified form;
\item it strictly increases progress and records;
\item it keeps the balance at or above the computed reserve.
\end{itemize}
\end{theorem}

Unbounded registers are essential here. No infinite run that strictly
progresses can fit into a finite full state, so the factor must be finite and
the counters must not. Solvency rests on the stated balance of credit and
debit. The greatest set is computed for the one fixed policy, not for every
solvent strategy. Disturbances are a unit class. This is a model-relative
result, not a theorem of physical persistence.

\subsection{Exact defects and robustness}

Required conditions cannot be averaged (Proposition~\ref{prop:scalar}). The
positive counterpart is an exact semantics of defects.

\begin{proposition}[Exact defect semantics]
% claim: TTS-02
\label{prop:defects}
Let each condition carry a labelled rational defect, and require defects and
tolerances to be nonnegative. A profile is accepted exactly when every required
defect is zero and every soft defect is at most its own tolerance. Acceptance is
downward closed, which means that lowering defects while keeping labels, types
and tolerances preserves it. Acceptance is not determined by the total defect sum, since two
profiles with the same schema and the same sum can differ.
\end{proposition}

\begin{proposition}[A bounded robustness box]
% claim: TTS-03
\label{prop:box}
In a rational-register extension of the two-copy model, on a fixed first event,
every correction inside a closed five-component box of half-width
\snapTtsZeroThreeHalfBoxRadius{} preserves strictly positive advances and both
complete handoff checks.
\end{proposition}

This is robustness of a stated, narrow kind, with one extension, one event and
one box.
It does not return to the original integer model, and it does not cover
arbitrary perturbations. A version stated for every perturbation of the
original model is false.
